\documentclass[12pt]{article}
\usepackage{amssymb,amsmath}
\usepackage[T2A]{fontenc}
\usepackage[utf8]{inputenc}
\usepackage[russian]{babel}
% \usepackage{latexsym}
% \RequirePackage[dvips]{graphicx}
% \usepackage{array,hhline}

% \usepackage{epsf,epic,eepic}
\input ris.tex
% \input cells.sty
\addrisoption{\risleftfalse}

\usepackage{pict2e}
\usepackage{tikz}

\DeclareRobustCommand{\No}{\ifmmode{\nfss@text{\textnumero}}\else\textnumero\fi}

%%
%% Три точки для обозначения делимости
%%
\catcode`\@=11
\def\kratno{\mathrel{\smash{\lower.5ex\hbox{$\,\vdots\,$}}}}
% Это прямые три точки для обозначения делимости
\def\nekratno{\mathrel{\mathpalette\c@ncel\kratno}}
\def\c@ncel#1#2{\m@th\ooalign{$\hfil#1\mkern1mu/\hfil$\crcr$#1#2$}}
% Это перечеркнутые три точки для обозначения неделимости
\def\divis{\smash{\lower.1ex\hbox{\,\hbox{.}\kern-.22em\lower-.6ex\hbox{.}
\kern-.52em\lower-1.2ex\hbox{.}\,}} }
% А это наклонные...

\hyphenation{чис-ло чис-ла чис-лу чис-лом чис-ле чис-лам чис-ла-ми чис-лах
че-ты-рех-уголь-ник че-ты-рех-уголь-ни-ка пря-мо-уголь-ник тре-уголь-ни-ка
тре-уголь-ни-ков вне-впи-сан-ной вне-впи-сан-ных}

\def\signed(#1){{\unskip\nobreak\hfil\penalty50\hskip1em\hbox{}\nobreak\hfil
                 \rm(\small\itshape #1\/\normalsize\rm)\parfillskip=0pt\finalhyphendemerits=0\par}}
        % macro for authorship's signature, from TeXbook
\def\be{\signed(С.~Берлов)}
\def\fp{\signed(Ф.~Петров)}
\def\frank{\signed(В.~Франк)}
\def\bah{\signed(Ф.~Бахарев)}
\def\guas{\signed(А.~Голованов)}
\def\as{\signed(А.~Смирнов)}
\def\dvk{\signed(Д.~Карпов)}
\def\dm{\signed(Д.~Максимов)}
\def\kk{\signed(К.~Кохась)}
\def\kuz{\signed(А.~Кузнецов)}
\def\hr{\signed(А.~Храбров)}
\def\ov{\signed(О.~Ванюшина)}
\def\io{\signed(О.~Иванова)}
\def\sop{\signed(Е.~Сопкина)}
\def\sviv{\signed(С.~Иванов)}
\def\suk{\signed(К.~Сухов)}
\def\dr{\signed(Д.~Ростовский)}
\def\pas{\signed(А.~Пастор)}
\def\azh{\signed(А.~Железняк)}
\def\vlas{\signed(Н.~Власова)}
\def\asol{\signed(А.~Солынин)}
\def\miv{\signed(М.~Иванов)}
\def\chuh{\signed(А.~Чухнов)}
\def\lev{\signed(Л.~Емельянов)}
\def\dsh{\signed(Д.~Ширяев)}
\def\zh{\signed(жюри)}
\def\bdzh{\signed(О.~Бадажкова)}

\sloppy
\pagestyle{empty}

\def\q#1.{{\bf #1. }}

%%
%% Макрокоманда для рисования клетчатых матриц с поправкой на две жирные линии
%%
\def\btable#1#2#3\endbtable{{%
\offinterlineskip\def\\{\cr\noalign{\hrule}}\halign
  {\strut\vrule\hbox to #1{\hfil#2\hfil}\vrule width 1.5 pt&&\hbox to #1{\hfil#2\hfil}\vrule\cr\noalign{\hrule}#3}}}

\begin{document}

\centerline{\sc САНКТ-ПЕТЕРБУРГСКАЯ}
\centerline{\sc ОЛИМПИАДА ШКОЛЬНИКОВ 2023 года ПО МАТЕМАТИКЕ}
\centerline{\sc I тур. \ 6 класс.}
\medskip
\hrule
\bigskip

\q1. Богатырь И.\,Муромец рубится с З.\,Горынычем, у которого 444 головы. Одним взмахом меча он может отсечь змею одну голову или сразу 10~голов (конечно, если их было не меньше 10).  Однако, если после отсечения осталось чётное число голов, то число голов сразу же удваивается. Приведите пример, как может действовать богатырь, чтобы отрубить З.\,Горынычу все его головы. 

\q2.  В ряд стоят 12 рыцарей и 12 лжецов (рыцари всегда говорят правду, а лжецы всегда лгут). Каждый произнёс либо фразу <<Слева от меня чётное число человек>>, либо фразу <<Слева от меня нечётное число человек>>. Могло ли оказаться, что всего было произнесено 17 фраз одного типа и 7 другого? Не забудьте обосновать свой ответ.

\q3. Костя задумал натуральное число $x$ и обнаружил, что некоторое четырехзначное число при делении на $x$ дает остаток $24$, а при делении на $x^2$ дает остаток 142. Найдите все возможные значения, которые может принимать $x$. Не забудьте обосновать свой ответ.

\q4. Клетки горизонтальной полоски $1\times 100$ пронумерованы слева направо числами от 1 до 100. В клетках с 1-й по 10-ю живет по одному кузнечику. Энтомолог Вася и кузнечики играют в игру: сначала ходит Вася, потом кузнечики, затем снова Вася, и т.\,д. За один ход Вася может перенести \emph{одного} из кузнечиков вправо не дальше чем на 6 клеток. На своем ходу \emph{все} кузнечики, находящиеся не в своем доме, перепрыгивают в соседнюю слева клетку. Докажите, что Вася не сможет собрать кузнечиков (в каком-то порядке) в клетках 50, 51, . . . , 59. В процессе игры разрешается нескольким кузнечикам находиться в одной клетке.


\clearpage

\centerline{\sc САНКТ-ПЕТЕРБУРГСКАЯ}
\centerline{\sc ОЛИМПИАДА ШКОЛЬНИКОВ 2023 года ПО МАТЕМАТИКЕ}
\centerline{\sc I тур. \ 7 класс.}
\medskip
\hrule
\bigskip

\q1. Богатырь И.\,Муромец рубится с З.\,Горынычем, у которого 444 головы. Одним взмахом меча он может отсечь змею одну голову или сразу 10~голов (конечно, если их было не меньше 10). Однако, если после отсечения осталось чётное число голов, то число голов сразу же удваивается. Сможет ли богатырь отрубить З.\,Горынычу все его головы?  Не забудьте обосновать свой ответ.

\q2. Каждый из 25 детей держит в руках табличку с ненулевым числом (возможно, отрицательным), все эти числа разные. Дети построились в~ряд по убыванию чисел (первое~--- самое большое), и Петя оказался \emph{десятым} по счету. Затем дети построились по убыванию чисел, обратных к исходным (напомним, что обратным к числу $a$ называется число $1/a$), и~Петя оказался \emph{шестнадцатым}. Наконец, дети построились по убыванию квадратов исходных чисел (все квадраты оказались разными). Каким по счету может оказаться Петя? Приведите все варианты и~объясните, почему других нет.

\q3. Натуральное число $N$ имеет больше 400 натуральных делителей (включая 1 и $N$). Все эти делители записали на доске. Саша стёр сто наибольших и сто наименьших из них. Среди оставшихся делителей оказалось поровну четных и нечетных. Докажите, что если бы вместо этого он стёр двести наибольших и двести наименьших делителей, среди оставшихся тоже оказалось бы поровну четных и нечетных.

\q4. В некоторых клетках доски $9 \times 9$ проведены диагонали (возможно, в некоторых клетках проведены сразу обе диагонали). Оказалось, что никакие четыре проведённые диагонали не имеют общего конца. Какое наибольшее число диагоналей может быть проведено?


\clearpage

\centerline{\sc САНКТ-ПЕТЕРБУРГСКАЯ}
\centerline{\sc ОЛИМПИАДА ШКОЛЬНИКОВ 2023 года ПО МАТЕМАТИКЕ}
\centerline{\sc I тур. \ 8 класс.}
\medskip
\hrule
\bigskip

\q1. На столе стоят 10 гирь разного веса. Известно, что сумма весов пяти самых лёгких гирь равна половине суммы весов остальных. Докажите, что сумма весов шести самых лёгких гирь меньше суммы весов остальных.

\q2. Каждый из 25 детей держит в руках табличку с ненулевым числом (возможно, отрицательным), все эти числа разные. Дети построились в~ряд по убыванию чисел (первое~--- самое большое), и Петя оказался \emph{десятым} по счету. Затем дети построились по убыванию чисел, обратных к исходным (напомним, что обратным к числу $a$ называется число $1/a$), и~Петя оказался \emph{шестнадцатым}. Наконец, дети построились по убыванию квадратов исходных чисел (все квадраты оказались разными). Каким по счету может оказаться Петя? Приведите все варианты и~объясните, почему других нет.

\q3. Внутри угла $ABC$ отмечена точка $D$. Известно, что $AB=2$, $CD=3$, $BC=4$, $\angle ABC=\angle BCD=60^\circ$. Точка $E$~--- середина отрезка $BD$. Найдите $AE$.

\q4. В некоторых клетках доски $2023 \times 2023$ проведены диагонали (возможно, в некоторых клетках проведены сразу обе диагонали). Оказалось, что никакие четыре проведённые диагонали не имеют общего конца. Какое наибольшее число диагоналей может быть проведено?

\q5. Натуральное число $N$ имеет больше 1200 натуральных делителей (включая 1 и $N$). Все эти делители записали на доске. Саша стёр 300 наибольших и 300 наименьших из них. Среди оставшихся делителей оказалось поровну четных и нечетных. Докажите, что среди всех делителей тоже поровну четных и нечетных.


\clearpage

\centerline{\sc САНКТ-ПЕТЕРБУРГСКАЯ}
\centerline{\sc ОЛИМПИАДА ШКОЛЬНИКОВ 2023 года ПО МАТЕМАТИКЕ}
\centerline{\sc I тур. \ 9 класс.}
\medskip
\hrule
\bigskip

\q1. Разность корней квадратного трёхчлена $P$ равна 4, а разность корней трёхчлена $P+6$ равна 8. Чему может быть равна разность корней квадратного трёхчлена $P+28$?

\q2. Каждый из 97 детей держит в руках табличку с ненулевым (положительным или отрицательным) числом, все эти числа разные. Дети построились в ряд по убыванию чисел (первое~--- самое большое), и~Петя оказался \emph{сороковым} по счету. Затем дети построились по убыванию чисел, обратных к исходным (напомним, что обратным к числу~$a$ называется число $1/a$), и Петя оказался \emph{шестидесятым}. Наконец, дети построились по убыванию квадратов исходных чисел (все квадраты оказались разными). Каким по счету может оказаться Петя? Приведите все варианты и объясните, почему других нет.

\q3. Точка $K$ лежит на биссектрисе $CL$ треугольника $ABC$. На стороне $AB$ выбрана такая точка $X$, что $AC=3KX=9AX$ и $KX\parallel AC$. Известно, что $AB=2$. Найдите периметр треугольника $ABC$.

\q4. Назовём натуральное число $n$ {\it полезным}, если у числа $n^2$ найдется такой делитель $d$, что число $n^2+101d$ является точным квадратом. Сколько полезных чисел содержится среди чисел от 1 до 1\,000\,000? 

\q5. Серёже на день рождения подарили доску, представляющую собой квадрат $150\times 150$ клеток, из которого вырезаны два квадрата $6\times 6$: содержащий левую верхнюю угловую клетку и содержащий правую верхнюю угловую клетку. В~каждой клетке этой доски он написал натуральное число от 1 до 8, в результате каждое число встречается на доске нечётное число раз. Уголок из трёх клеток назовём {\it удачным}, если все три числа в его клетках равны, или они все различны. Сережа подсчитал количество удачных уголков на доске. Какое наименьшее количество у~него могло получиться?


\clearpage

\centerline{\sc САНКТ-ПЕТЕРБУРГСКАЯ}
\centerline{\sc ОЛИМПИАДА ШКОЛЬНИКОВ 2023 года ПО МАТЕМАТИКЕ}
\centerline{\sc I тур. \ 10 класс.}
\medskip
\hrule
\bigskip

\q1. Разность корней квадратного трёхчлена $P$ равна $2$, а  разность корней трёхчлена $P + 3$ равна $4$. Чему может быть равна разность корней квадратного трёхчлена $P + 8$?

\q2. В каждой клетке прямоугольной таблицы стоит 0 или 1. Известно, что в каждой строке таблицы стоит не меньше 24 единиц, а в каждом столбце~--- не меньше 24 нулей. Какое наименьшее количество клеток может быть в таблице?

\q3. В лагере отдыхало несколько школьников, один из них~--- Вася. У~каждого школьника в лагере было ровно 19 друзей. Через несколько дней Вася и его 19 друзей уехали, в результате чего у~всех оставшихся в лагере стало по 17 друзей. Какое наибольшее число школьников могло остаться в лагере?

\q4. Натуральные числа $a$, $b$, $c$ таковы, что $b^2a - 1$ делится на $ac - 1$, $a^2c - 1$ делится на $bc - 1$ и $c^2b - 1$ делится на $ab - 1$. Докажите, что хотя бы одно из чисел $a$, $b$, $c$~--- точный квадрат.

\q5. В выпуклом 100-угольнике $M$ все углы равны. При повороте на некоторый угол из него получился 100-угольник $N$. Стороны $M$ покрасили в зеленый цвет, а стороны $N$~--- в черный. Пересечение многоугольников $M$ и $N$ оказалось 200-угольником, в котором зеленые и черные стороны чередуются. Сумма квадратов длин чёрных сторон равна 100. Чему может быть равна сумма квадратов длин зеленых сторон?


\clearpage

\centerline{\sc САНКТ-ПЕТЕРБУРГСКАЯ}
\centerline{\sc ОЛИМПИАДА ШКОЛЬНИКОВ 2023 года ПО МАТЕМАТИКЕ}
\centerline{\sc I тур. \ 11 класс.}
\medskip
\hrule
\bigskip

\q1. В каждой клетке таблицы $m\times n$ стоит крестик или нолик. При этом в каждой строке стоит не меньше 19 крестиков, а в каждом столбце~--- не меньше 19 ноликов. Докажите, что $m+n\geq 76$.

\q2. Даны квадратные трехчлены $f(x)$ и $g(x)$ со старшими коэффициентами, равными~1. Наименьшее значение трёхчлена $f(x)$ равно~10, наименьшее значение $g(x)$ равно 20. А наименьшее значение квадратного трёхчлена $f(x)+g(x)$ равно 50 и достигается в точке $x=x_0$. Найдите $f(x_0)$ и $g(x_0)$.

\q3. Функции $p$ и $q$ определены на множестве натуральных чисел, не превосходящих 400, и принимают натуральные значения, не превосходящие 200. Докажите, что уравнение 
$$
p(x)-p(y)=q(x)-q(y)
$$ 
имеет решение в натуральных числах $x\ne y$.

\q4. Через вершины $A$ и $B$ равнобедренного треугольника $ABC$ ($AB=BC$) проведена окружность, пересекающая отрезки $AC$ и~$BC$ в точках $X$ и $Y$ соответственно. Описанная окружность треугольника $CXY$ пересекает отрезок $BX$ в точке $Z$. Оказалось, что $AB=2CZ$. Найдите $BZ/ZX$.

\q5. Натуральное число $N$ имеет больше 1000 натуральных делителей (включая 1 и $N$). Все эти делители выписаны на доске. Саша стёр 250 наибольших и 250 наименьших из них. Среди оставшихся делителей оказалось поровну делящихся на~13 и не делящихся на~13. Докажите, что среди всех делителей тоже поровну делящихся и не делящихся на~13.

\end{document}
